% Generated from locale/tr/content/first-order-logic/sequent-calculus/derivations.tex; do not hand-edit.


\iftag{FOL}
      {\olfileid[tr]{fol}{seq}{der}}
      {\olfileid[tr]{pl}{seq}{der}}

\olsection{\usetoken{P}{derivation}}

\begin{explain}
Başlangıç sekantlarının biçimini belirttik ve çıkarım kurallarını verdik.
Sekant hesabındaki !!^{derivation}s bunlardan tümevarımsal olarak üretilir:
her !!{derivation} ya tek başına bir başlangıç sekantıdır ya da bir çıkarımın
izlediği !!{derivation}s arasından bir veya ikisini içerir.
\end{explain}

\begin{defn}[$\Log{LK}$-!!{derivation}]
Bir $S$ sekantının \emph{$\Log{LK}$-!!{derivation}{}i}, aşağıdaki koşulları
sağlayan sonlu bir sekant ağacıdır:
\begin{enumerate}
\item Ağacın en üstteki sekantları başlangıç sekantlarıdır.
\item Ağacın en alttaki sekantı $S$'dir.
\item Ağaçta $S$ dışında kalan her sekant, sonucu o sekantın hemen altında
  duran bir çıkarım kuralının doğru bir uygulamasının öncülüdür.
\end{enumerate}
Bu durumda $S$'ye !!{derivation}{}in \emph{son sekantı}, $S$'ye de
\emph{$\Log{LK}$'de !!{derivable}} (veya $\Log{LK}$-!!{derivable}) deriz.
\end{defn}

\begin{ex}
Her başlangıç sekantı, örneğin $!C \Sequent !C$, !!a{derivation}{}dir. Bu
!!{derivation}{}den, örneğin $\LeftR{\Weakening}$ kuralını uygulayarak yeni bir türetim
elde edebiliriz:
\begin{prooftree}
\Axiom$ \Gamma \fCenter \Delta $
\RightLabel{\LeftR{\Weakening}}
\UnaryInf$ !A, \Gamma \fCenter \Delta$
\end{prooftree}
Ancak kural genel anlaşılmalıdır: kuraldaki $!A$ yerine herhangi bir
!!{sentence}, örneğin $!D$ de koyabiliriz. Öncül, başlangıç sekantımız
$!C \Sequent !C$ ile eşleşiyorsa hem $\Gamma$ hem de $\Delta$ yalnızca
$!C$'dir ve sonuç $!D,!C \Sequent !C$ olur. Dolayısıyla aşağıdaki bir
!!a{derivation}{}dir:
\begin{prooftree}
\Axiom$ !C \fCenter !C $
\RightLabel{\LeftR{\Weakening}}
\UnaryInf$ !D, !C \fCenter !C$
\end{prooftree}
Şimdi solda duran !!{sentence}s arasından ikisinin yerini değiştirmemizi sağlayan başka bir
kuralı, örneğin $\LeftR{\Exchange}$ kuralını uygulayabiliriz. Buna göre
aşağıdaki de doğru bir !!{derivation}{}dir:
\begin{prooftree}
\Axiom$ !C \fCenter !C $
\RightLabel{\LeftR{\Weakening}}
\UnaryInf$ !D, !C \fCenter !C$
\RightLabel{\LeftR{\Exchange}}
\UnaryInf$ !C, !D \fCenter !C$
\end{prooftree}
Şu biçimde verilen kuralın bu uygulamasında
\begin{prooftree}
\Axiom$ \Gamma, !A, !B, \Pi \fCenter \Delta $
\RightLabel{\LeftR{\Exchange}}
\UnaryInf$ \Gamma, !B, !A, \Pi \fCenter \Delta,$
\end{prooftree}
$\Gamma$ ile $\Pi$ boş, $\Delta$ ise $!C$'dir; $!A$ ile $!B$ rollerini de
sırasıyla $!D$ ile $!C$ oynar. Hemen hemen aynı biçimde,
\begin{prooftree}
\Axiom$ !D \fCenter !D $
\RightLabel{\LeftR{\Weakening}}
\UnaryInf$ !C, !D \fCenter !D$
\end{prooftree}
ifadesinin de !!a{derivation} olduğunu görürüz. Artık !!{derivation}s elimizde
olduğuna göre bu ikisini alıp $\RightR{\land}$ ile birleştirebiliriz. Bu kural şuydu:
\begin{prooftree}
\Axiom$\Gamma \fCenter \Delta, !A$
\Axiom$ \Gamma \fCenter \Delta, !B$
\RightLabel{\RightR{\land}}
\BinaryInf$ \Gamma \fCenter \Delta, !A \land !B $
\end{prooftree}
Bizim durumumuzda öncüller, söz konusu !!{derivation}s için son sekantlardır
ve kuralın öncülleriyle eşleşmelidir. Bu, $\Gamma$'nın $!C,!D$, $\Delta$'nın boş,
$!A$'nın $!C$ ve $!B$'nin $!D$ olduğu anlamına gelir. Dolayısıyla çıkarım
doğru olacaksa sonuç $!C,!D \Sequent !C \land !D$ olur.
\begin{prooftree}
\Axiom$ !C \fCenter !C $
\RightLabel{\LeftR{\Weakening}}
\UnaryInf$ !D, !C \fCenter !C$
\RightLabel{\LeftR{\Exchange}}
\UnaryInf$ !C, !D \fCenter !C$
\Axiom$ !D \fCenter !D $
\RightLabel{\LeftR{\Weakening}}
\UnaryInf$ !C, !D \fCenter !D$
\RightLabel{\RightR{\land}}
\BinaryInf$ !C, !D \fCenter !C \land !D $
\end{prooftree}
Elbette öncüllerin sırasını tersine de çevirebiliriz; o zaman $!A$, $!D$ ve
$!B$, $!C$ olur.
\begin{prooftree}
\Axiom$ !D \fCenter !D $
\RightLabel{\LeftR{\Weakening}}
\UnaryInf$ !C, !D \fCenter !D$
\Axiom$ !C \fCenter !C $
\RightLabel{\LeftR{\Weakening}}
\UnaryInf$ !D, !C \fCenter !C$
\RightLabel{\LeftR{\Exchange}}
\UnaryInf$ !C, !D \fCenter !C$
\RightLabel{\RightR{\land}}
\BinaryInf$ !C, !D \fCenter !D \land !C $
\end{prooftree}
\end{ex}

